\scanpage{27}
\mspage{14}

% top centre, above proof
\margin{\unsure{OPS|OBS} $k$ infini}

% upper third
\uline{Dém} Un $(\lambda_i)\in N(u_*)$ tel que les comp. sont $\ne0$ est connu \unsure{quand} on \unsure{connait|connaît} trois comp. consécutives, disons $\lambda_1,\lambda_2,\lambda_3$ (on utilise le fait que les $u_i$ sont $\ne0$), \deleted{et on en conclut} i.e.

% upper middle, two-line map
\diagram{\begin{tikzcd}[ampersand replacement=\&,row sep=2pt,column sep=large]
N(u_*)\arrow[r,dash] \& k^3\\
(\lambda_i)\arrow[r,mapsto] \&(\lambda_1,\lambda_2,\lambda_3)
\end{tikzcd}}

% right of map
a un noyau

% left margin rising to the right
% slanted
\margin{Lemme (si $N(u_*)\not\subset\bigcup H_i$)}

% middle, heavily reworked sentence
\uline{\unsure{V|v}}\quad \inserted{\unsure{dont de quelque translaté|dont des translatés} \deleted{de} $N(u_*)$} \unsure{contenu|contenus} \deleted{dans} la réunion des hyperplans $H_i$ \deleted{d'équations} \inserted{\unsure{séparés|sép.} par un pt} \unsure{de coordonnées|d'ordonnées|d'équations} $\lambda_i=0$ $(i\in\mathbb Z/n\mathbb Z)$ -- donc, $k$ étant infini, \deleted{\unsure{l'on a bien}} ce noyau est réduit à un pt, i.e. est nul. Donc

% centre, injectivity and bounds
$N(u_*)\to k^3$ est \uline{injectif}. Donc $1\leq\dim N(u_*)\leq3$. Mais \deleted{\unsure{OPS} \unsure{quelques} $u_i$}

% lower half
On a $\dim N(u_*)\ne3$, i.e. $N(u_*)\mathbin{\unsure{-|\to}}k^3$ n'est pas surjectif -- \deleted{il} \inserted{l'image} ne contient pas p.ex $(1,0,0)$ (car $\lambda_1u_0\deleted{\unsure{+}}-\lambda_2(u_1+u_2)+\lambda_3u_3$ implique $u_0=0$, \unsure{contraire|contrmt} à l'hyp. On veut regarder le cas $\dim N(u_*)=2$. On

% lower third, image and kernel
{\small
\[
\begin{aligned}
\deleted{N(u_*)\longrightarrow}\quad N(u_*)&\simeq\operatorname{Im}(N(u_*)\to k^3)
 \subset\{(\lambda_1,\lambda_2,\lambda_3)\mid\lambda_1u_0-\lambda_2(u_1+u_2)\\
&\hspace{23em}+\lambda_3u_3=0\}\\
\deleted{\text{donc}}\quad &=\operatorname{Ker}\big((\lambda_1,\lambda_2,\lambda_2)\mathbin{\unsure{-|\to|\mapsto}}\lambda_1u_0-\lambda_2(u_1+u_2)+\lambda_3u_3\big)
\end{aligned}
\]
}

% bottom, dimensions and two cases
\[
\begin{aligned}
\text{donc}\quad\dim N(u_*)&\leq\dim\operatorname{Ker}=\deleted{\unsure{f}}3-
\deleted{\unsure{\dim}}\operatorname{rg}(u_0,u_1+u_2,u_3)\\
&=\begin{cases}1&\text{si }\mathrm{rg}=2\\2&\text{si }\mathrm{rg}=1\end{cases}
\end{aligned}
\]

% archive sheet number: 33
